\section{সমাজজীবনে তথ্য ও যোগাযোগ প্রযুক্তির প্রভাব}
\label{sec:ch1-social}
\subsection{সুযোগ ও সামাজিক পরিবর্তন}
মোবাইল, ই-মেইল, তাৎক্ষণিক বার্তা, ভিডিও কল ও সামাজিক যোগাযোগে পরিবার, শিক্ষাপ্রতিষ্ঠান, অফিস এবং আন্তর্জাতিক যোগাযোগের সময় ও ব্যয় কমে। অনলাইন পাঠ, ডিজিটাল গ্রন্থাগার, ভিডিও ও পরীক্ষা জ্ঞানচর্চায় প্রবেশের সুযোগ বাড়ায়। দূরবর্তী চিকিৎসা, অ্যাপয়েন্টমেন্ট, রিপোর্ট, স্বাস্থ্য অ্যাপ ও জরুরি যোগাযোগে প্রত্যন্ত অঞ্চলের মানুষও পরামর্শ পেতে পারেন। সেবাগুলোর কার্যপদ্ধতি বিশ্বগ্রামের প্রয়োগ অংশে আছে।
দূরের কাজ, স্বাধীন চুক্তির কাজ ও ডিজিটাল উদ্যোগ আয়ের নতুন ক্ষেত্র তৈরি করে। পুনরাবৃত্ত কাজ কমলে রক্ষণাবেক্ষণ, তথ্য বিশ্লেষণ ও সেবার নকশার মতো দক্ষতা দরকার হয়। পুরোনো দক্ষতার কর্মীর পুনঃপ্রশিক্ষণ ও ন্যায্য সুযোগ জরুরি; সব পেশা একসঙ্গে বিলুপ্ত হবে এমন সিদ্ধান্ত ঠিক নয়।
সংগীত, ভিডিও, গেম ও সামাজিক প্ল্যাটফর্ম বিনোদন এবং জনসচেতনতার সুযোগ বাড়ায়। স্থানীয় ভাষা, সাহিত্য ও শিল্প বিশ্বে পৌঁছাতে পারে; অন্য ভাষা ও সংস্কৃতিও শেখা যায়। নিজের সংস্কৃতির ডিজিটাল সংগ্রহশালা ও মানসম্মত বিষয়বস্তু তৈরি উপযোগী প্রয়োগ।
\subsection{বৈষম্য, সংস্কৃতি ও সীমাবদ্ধতা}
সেবা অনলাইন হলেই সবার সুযোগ সমান হয় না। ডিজিটাল বৈষম্যের মূল ধারণা \ref{sec:ch1-digital-divide} অংশে আছে। অনলাইন ক্লাসে কম ক্ষমতার ফোন ও সীমিত ডেটার শিক্ষার্থীর জন্য ছোট ফাইল, ডাউনলোডযোগ্য পাঠ, ক্যাপশন, বিকল্প মুদ্রিত পাঠ, প্রশিক্ষণ ও ভাগাভাগি করে ব্যবহারের সুযোগ দরকার।
ভুল তথ্য, গোপনীয়তার ক্ষতি, নিম্নমানের পাঠ, মনোযোগের ব্যাঘাত ও অন্যের কাজ নিজের বলে চালানো শিক্ষার সুফল কমায়। স্থানীয় ভাষার অবহেলা, একরকম সংস্কৃতির চাপ বা ভুল উপস্থাপনেও ক্ষতি হতে পারে। স্বাস্থ্যতথ্যের সুরক্ষা এবং অনলাইনে হয়রানির প্রতিকার দরকার; সংশ্লিষ্ট নৈতিকতা \ref{sec:ch1-ethics} অংশে আছে।
\subsection{স্বাস্থ্য ও সুস্থ অভ্যাস}
দীর্ঘক্ষণ একই ভঙ্গিতে বসা, অনুপযুক্ত আলো ও অতিরিক্ত পর্দা ব্যবহার চোখের অস্বস্তি, ঘুমের সমস্যা ও শারীরিক চাপ বাড়াতে পারে। নিয়মিত বিরতি, উপযুক্ত বসার ব্যবস্থা ও সময় নির্ধারণ দরকার। সরাসরি সামাজিক সম্পর্ক বজায় রাখতে হয়। অতিরিক্ত গেম বা অনলাইন ব্যবহারে আসক্তি কিংবা হয়রানির লক্ষণ থাকলে বিশ্বস্ত ব্যক্তি বা পেশাজীবীর সহায়তা নেওয়া উচিত।
\begin{practicebox}
\begin{enumerate}
\item সমাজজীবনে ICT-এর তিনটি ইতিবাচক প্রভাব লেখ।
\item ICT-এর অতিরিক্ত ব্যবহারের দুটি ক্ষতি লেখ।
\item Telemedicine কীভাবে সমাজে প্রভাব ফেলে?
\end{enumerate}
\end{practicebox}
